\section{Experimental Setup}\label{sec:setup}
\textbf{Hardware.} \TBD{Raspberry Pi~5 (RAM), OS and kernel, active/passive
cooler, enclosure, power supply, ambient temperature range, camera/source}.
A second device (\TBD{optional: Jetson / Pi~4 / x86}) would test generality.

\textbf{Software.} \shade{} commit \TBD{hash}; ONNX Runtime~\cite{onnxruntime}
\TBD{version}; models: YOLOv8n~\cite{ultralytics2023} at $640$ and $320$ input,
and an INT8 variant produced with static QDQ quantisation~\cite{jacob2018quant}
\TBD{calibration set and size}. The Detect head is kept in FP32 because
quantising it zeroed class scores in our tests.

\textbf{Ladder.} Table~\ref{tab:rungs} lists the rungs; operating points are
measured, not assumed (E1). Pareto-dominated rungs are pruned with the provided
script.

\textbf{Workloads.} \TBD{video clips: description, resolution, duration,
fraction of static time, availability and licence}. Ground truth
\TBD{source of labels, or agreement against the nominal rung}.

\textbf{Arms.} Static nominal, static emergency, \shade{} (default policy),
\shade{} with gate/shift, and \shade{} with the forecast advisor. Each arm runs
\TBD{$\ge5$} times for \TBD{30} minutes in the same enclosure at the same ambient
temperature, in randomised order, with a cool-down to baseline between runs.

\textbf{Metrics.} Mean and 5th-percentile output frame rate; share of time
firmware-throttled and above $70^\circ$C; maximum temperature; rung changes per
hour; agreement F1 against the nominal rung (IoU~0.5, and 0.75 for scheduling);
detector calls per frame; swap stall from the audit log.
